\subsection{正线性型的延拓}

命题 1. --- 设 E 是预序向量空间，V 是 E 的一个向量子空间，使得 E 的每个元素都被 V 的某个元素上有界。给定 V 上关于 V 的预序向量空间结构（由 E 的结构诱导）为正的线性型 f，存在 E 上正线性型组成的一个非空集 \(S_{f}\)，其中每个线性型都是 f 的延拓。若 \(h \in S_{f}\)，则诸值 \(h(a)\)（对 \(a \in E\)）都位于区间 \([\alpha', \alpha'']\) 中，其中

\[
\alpha^ {\prime} = \sup _ {z \in \mathbf {V}, z \leqslant a} f (z), \quad \alpha^ {\prime \prime} = \inf _ {y \in \mathbf {V}, y \geqslant a} f (y).\tag{1}
\]

I. 特殊情形。

首先设 \(E = V + R a\)。由于当 \(a \in V\) 时命题是平凡的，我们限于 \(a \notin V\) 的情形。对 V 的假设蕴含集 \(A''\)（诸 \(y \in V\)（满足 \(a \leqslant y\)）组成的）非空；类似地，集 \(A'\)（诸 \(z \in V\)（满足 \(-z \geqslant -a\)，即 \(z \leqslant a\)）组成的）非空。对 \(y \in A''\) 及 \(z \in A'\)，我们有 \(z \leqslant a \leqslant y\)，从而由假设 \(f(z) \leqslant f(y)\)。于是 \(\alpha', \alpha''\) 有限且 \(\alpha' \leqslant \alpha''\)。E 上延拓 f 的任何线性型 \(f_1\) 都由 \(f_1(a)\) 完全确定，且对所有 \(\lambda \in R\) 及所有 \(x \in V\)，我们有

\[
f _ {1} (x + \lambda a) = f (x) + \lambda f _ {1} (a).
\]

于是 \(f_{1}\) 是正的，当且仅当关系

\[
x \in \mathbf {V}, \quad \lambda \in \mathbf {R}, \quad x + \lambda a \geqslant 0\tag{2}
\]

蕴含

\[
f (x) + \lambda f _ {1} (a) \geqslant 0.\tag{3}
\]

由于 \(f(\mu x) = \mu f(x)\) 且关系 \(x \geqslant 0\) 与 \(\mu x \geqslant 0\)（对 \(\mu > 0\)）等价，只需证明 (2) 在特殊情形 \(\lambda = 0, \lambda = 1\) 与 \(\lambda = -1\) 下蕴含 (3)。对 \(\lambda = 0\)，(2) 蕴含 (3) 这一事实由 \(f\) 是正的这一假设得出。对 \(\lambda = 1\)，说 (2) 蕴含 (3) 意味着：对 \(-x \in A'\)，有 \(f_1(a) \geqslant f(-x)\)，即 \(f_1(a) \geqslant \alpha'\)；对 \(\lambda = -1\)，(2) 蕴含 (3) 意味着：对 \(x \in A''\)，有 \(f(x) \geqslant f_1(a)\)，即 \(f_1(a) \leqslant \alpha''\)。因此命题在这种情形下得证。

\paragraph{II. 一般情形。}

设 \(\mathfrak{F}\) 是由对 \((\mathrm{W}, g)\) 组成的集，其中 \(\mathrm{W}\) 是 \(\mathrm{E}\) 的包含 \(\mathrm{V}\) 的向量子空间，而 \(g\) 是 \(\mathrm{W}\) 上的正线性型且为 \(f\) 的延拓。我们给 \(\mathfrak{F}\) 排序：置 \((\mathrm{W}, g) \leqslant (\mathrm{W}', g')\)（若 \(\mathrm{W} \subset \mathrm{W}'\) 且 \(g'\) 是 \(g\) 的延拓）。显然 \(\mathfrak{F}\) 是归纳的，且由 S, III, § 2.4 的 th. 2，存在一个极大元素 \((\mathrm{W}_0, g_0)\)。设 \(\mathrm{W}_0 \neq \mathrm{E}\)。则存在向量 \(b \notin \mathrm{W}_0\)，且若 \(\mathrm{W}_1 = \mathrm{W}_0 + \mathbf{R}b\)，则上面的特殊情形表明存在 \(\mathrm{W}_1\) 上的正线性型且为 \(g_0\) 的延拓；这与 \((\mathrm{W}_0, g_0)\) 极大这一假设矛盾。于是 \(\mathrm{W}_0 = \mathrm{E}\)，命题的第一部分得证。当 \(a \in \mathrm{V}\) 时，第二个断言在 \(\alpha' = \alpha'' = f(a)\) 时显然成立；反之，若 \(a \notin \mathrm{V}\) 且置 \(\mathrm{V}_1 = \mathrm{V} + \mathbf{R}a\)，则第二个断言由证明中的特殊情形 I 得出。

推论. --- 在赋有相容预序结构的拓扑向量空间 E 中，设 P 是 E 中由 \(\geqslant 0\) 的元素组成的集。设 V 是 E 的一个向量子空间，至少包含 P 的一个内点 \(x_{0}\)。则 V 上每个正线性型都可以延拓为 E 上的正线性型。

由 prop. 1，只需证明对每个 \(x \in \mathbf{E}\)，存在 \(x' \in \mathbf{V}\) 使得 \(x' - x \in \mathbf{P}\)。现设 \(\mathbf{U}\) 是 \(\mathbf{E}\) 中 0 的一个邻域，使得 \(x_0 + \mathbf{U} \subset \mathbf{P}\)。则 \(x + x_0 + \mathbf{U} \subset x + \mathbf{P}\)，因此存在 \(\varepsilon\) 满足 \(0 < \varepsilon < 1\) 且点 \(y = x_0 + (1 - \varepsilon)x\) 属于 \(x + \mathbf{P}\)；于是每个形如 \(x + \lambda(y - x)\) 的点都属于 \(x + \mathbf{P}\)（对 \(\lambda > 0\)）。若取 \(\lambda = 1 / \varepsilon\)，则 \(x + \lambda(y - x) = \lambda x_0 \in \mathbf{V}\)，由此得出结论。

若不假设 V 包含 P 的一个内点，则推论的结论不一定成立，即使 E 是有限维的且 \(P \cap V\) 包含在 V 中为内点的点（II, 第 91 页, exerc. 25, b)）。
% semantic-fragments: legacy-input-seam
\relax{}
